\newpage
\Article{61}
Was in Artikel~2 der vorigen Abhandlung über den Begriff des biquadratischen Restes und Nichtrestes gesagt wurde, dehnen wir auch auf die Arithmetik komplexer Zahlen aus. Wie dort beschränken wir auch hier die Untersuchung auf Moduln, die Primzahlen sind. Zugleich ist meist stillschweigend vorauszusetzen, dass der Modul so gewählt ist, dass er unter seinen Assoziierten primär, also $\equiv1$ nach dem Modul $2+2i$ ist, und dass die Zahlen, deren Charakter als biquadratische Reste oder Nichtreste betrachtet wird, durch den Modul nicht teilbar sind.

Für einen gegebenen Modul könnten somit die durch ihn nicht teilbaren Zahlen in drei Klassen eingeteilt werden: Die erste enthielte die biquadratischen Reste, die zweite diejenigen biquadratischen Nichtreste, die quadratische Reste sind, die dritte die quadratischen Nichtreste.

\SourceLeaf{141}%
Doch auch hier ist es besser, anstelle der dritten Klasse zwei aufzustellen, sodass es insgesamt vier gibt.

Wird eine beliebige primitive Wurzel als Basis genommen, so haben die biquadratischen Reste durch $4$ teilbare Indizes, also solche der Form $4n$; diejenigen Nichtreste, die quadratische Reste sind, haben Indizes der Form $4n+2$; schließlich haben die quadratischen Nichtreste teils Indizes der Form $4n+1$, teils solche der Form $4n+3$. Auf diese Weise ergeben sich zwar vier Klassen, doch wäre der Unterschied zwischen den beiden letzten kein absoluter: Er hinge von der Wahl der als Basis genommenen primitiven Wurzel ab. Denn man sieht leicht, dass die eine Hälfte der primitiven Wurzeln einem gegebenen quadratischen Nichtrest einen Index der Form $4n+1$, die andere Hälfte einen Index der Form $4n+3$ gibt. Um diese Mehrdeutigkeit zu beseitigen, setzen wir stets eine primitive Wurzel voraus, bei der der Index $\frac14(p-1)$ zur Zahl $+i$ gehört (vgl. Artikel~55, VI). So entsteht eine Klasseneinteilung, die sich kürzer und unabhängig von primitiven Wurzeln folgendermaßen aussprechen lässt.

\Needspace{10\baselineskip}
Die \emph{erste} Klasse enthalte diejenigen Zahlen $k$, für die $k^{\frac14(p-1)}\equiv1$ gilt; diese Zahlen sind die biquadratischen Reste des Moduls.

Die \emph{zweite} Klasse enthalte diejenigen, für die $k^{\frac14(p-1)}\equiv i$ gilt.

Die \emph{dritte} diejenigen, für die $k^{\frac14(p-1)}\equiv-1$ gilt.

Schließlich die \emph{vierte} diejenigen, für die $k^{\frac14(p-1)}\equiv-i$ gilt.

Die dritte Klasse enthält diejenigen biquadratischen Nichtreste, die quadratische Reste sind; die quadratischen Nichtreste verteilen sich auf die zweite und vierte Klasse.

Den Zahlen dieser Klassen legen wir der Reihe nach die \emph{biquadratischen Charaktere} $0,1,2,3$ bei. Definieren wir den Charakter $\lambda$ einer Zahl $k$ nach dem Modul $m$ als den Exponenten jener Potenz von $i$, der die Zahl $k^{\frac14(p-1)}$ kongruent ist, so sind offenbar nach dem Modul $4$ kongruente Charaktere als gleichwertig anzusehen. Vorerst ist dieser Begriff auf Moduln beschränkt, die Primzahlen sind: In der Fortsetzung dieser Untersuchungen werden wir zeigen, wie er auch auf zusammengesetzte Moduln angewandt werden kann.

\Needspace{10\baselineskip}
\Article{62}
Um eine ausgedehnte Untersuchung von Zahlenbeispielen zu den Charakteren der Zahlen zu erleichtern, fügen wir hier eine übersichtliche Tabelle hinzu. Mit ihrer Hilfe lässt sich der Charakter jeder gegebenen Zahl bezüglich eines Moduls, dessen Norm den Wert $157$ nicht übersteigt, mit wenig Arbeit erhalten, sofern die folgenden Bemerkungen beachtet werden.

\SourceLeaf{142}%
Da der Charakter einer zusammengesetzten Zahl gleich der Summe der Charaktere ihrer einzelnen Faktoren oder ihr nach dem Modul $4$ kongruent ist, genügt es, wenn wir für einen gegebenen Modul die Charaktere der Primzahlen angeben können. Da ferner die Charaktere der Einheiten $-1,i,-i$ offenbar den Zahlen $\frac12(p-1),\frac14(p-1),\frac34(p-1)$ nach dem Modul $4$ kongruent sind, genügt es auch, die Charaktere der unter ihren Assoziierten primären Zahlen anzugeben. Da schließlich nach dem Modul $m$ kongruente Moduln denselben Charakter haben, genügt es, die Charaktere solcher Zahlen in die Tabelle aufzunehmen, die in einem System absolut kleinster Reste enthalten sind. Außerdem zeigt eine ähnliche Schlussfolgerung wie in Artikel~58, dass, wenn für den Modul $a+bi$ der Charakter der Zahl $A+Bi$ gleich $\lambda$ ist und für den Modul $a-bi$ $\lambda'$ der Charakter der Zahl $A-Bi$ ist, stets $\lambda\equiv-\lambda'\pmod4$ gilt, oder gleichbedeutend $\lambda+\lambda'$ durch $4$ teilbar ist. Daher genügt es, Moduln in die Tabelle aufzunehmen, in denen $b$ entweder $0$ oder positiv ist.

\Needspace{16\baselineskip}
Wird zum Beispiel der Charakter der Zahl $11-6i$ bezüglich des Moduls $-5-6i$ gesucht, so ersetzen wir diese Zahlen durch $11+6i,-5+6i$. Sodann bestimmen wir (Artikel~43) den absolut kleinsten Rest der Zahl $11+6i$ nach dem Modul $-5+6i$; er ist $-1-4i=-1\times(1+4i)$. Da für den Modul $-5+6i$ der Charakter von $-1$ gleich $30$, der Charakter der Zahl $1+4i$ nach der Tabelle dagegen $2$ ist, ergibt sich $32$ oder $0$ als Charakter der Zahl $11+6i$ für den Modul $-5+6i$ und nach der letzten Bemerkung somit auch als Charakter der Zahl $11-6i$ für den Modul $-5-6i$. Wird ebenso der Charakter der Zahl $-5+6i$ bezüglich des Moduls $11+6i$ gesucht, so zerlegt sich ihr absolut kleinster Rest $1-5i$ in die Faktoren $-i\;1+i,\;3-2i$, denen die Charaktere $117,0,1$ entsprechen. Der gesuchte Charakter ist daher $118$ oder $2$; derselbe Charakter ist auch der Zahl $-5-6i$ bezüglich des Moduls $11-6i$ beizulegen.

\[
\begin{array}{r|c|l}
\text{Modul.}&\text{Charakter.}&\text{Zahlen.}\\\hline
-3&3&1+i\\
+3+2i&3&1+i\\
+1+4i&1&-1+2i\\
&3&1+i\\
-5+2i&0&-1-2i\\
&1&1+i\\
&2&-1+2i\\
-1+6i&0&-3\\
&1&1+i,\ -1+2i
\end{array}
\]

\SourceLeaf{143}%
\[
\begin{array}{r|c|l}
\text{Modul.}&\text{Charakter.}&\text{Zahlen.}\\\hline
-1+6i&2&-1-2i\\
+5+4i&0&1+i\\
&1&-3\\
&3&-1+2i,\ -1-2i\\
-7&0&-3\\
&1&-1+2i,\ 3-2i\\
&2&1+i\\
&3&-1-2i,\ 3+2i\\
+7+2i&0&1+i,\ 3+2i,\ 3-2i,\ 1-4i\\
&1&-3\\
&2&-1-2i,\ 1+4i\\
&3&-1+2i\\
-5+6i&0&1+i,\ -3,\ 3+2i,\ 3-2i\\
&1&1-4i\\
&2&1+4i\\
&3&-1+2i,\ -1-2i\\
-3+8i&0&-1+2i,\ 3-2i,\ 1-4i\\
&1&1+i,\ 3+2i\\
&2&-3\\
&3&-1-2i,\ 1+4i,\ -5+2i\\
+5+8i&0&-1-2i\\
&1&-5-2i,\ -1+6i\\
&2&-1+2i,\ 3-2i\\
&3&1+i,\ -3,\ 3+2i,\ 1+4i,\ 1-4i\\
+9+4i&0&-1+2i,\ 3+2i\\
&1&1+i,\ -1-2i,\ 3-2i\\
&2&-3,\ 1+4i\\
&3&1-4i,\ -5+2i\\
-1+10i&0&1+i,\ -1+2i,\ -1-2i,\ 3+2i\\
&1&-3\\
&2&3-2i,\ -5+2i,\ 5-4i\\
&3&1+4i,\ 1-4i
\end{array}
\]

\SourceLeaf{144}%
\[
\begin{array}{r|c|l}
\text{Modul.}&\text{Charakter.}&\text{Zahlen.}\\\hline
+3+10i&1&1+i,\ -1-2i,\ 1-4i\\
&2&-3,\ 3+2i,\ 1+4i,\ -5-2i\\
&3&-1+2i,\ 3-2i\\
-7+8i&0&1+i,\ -7\\
&1&3+2i,\ 3-2i,\ 1-4i,\ -5-2i\\
&2&-1-2i,\ 1+4i,\ -5+2i,\ -1-6i\\
&3&-1+2i,\ -3,\ -1+6i\\
-11&0&-3\\
&1&1+i,\ 3-2i,\ 1+4i,\ -5+2i,\ 5+4i\\
&2&-1+2i,\ -1-2i\\
&3&3+2i,\ 1-4i,\ -5-2i,\ 5-4i\\
-11+4i&0&1+i,\ -1+2i,\ 3+2i,\ 5+4i\\
&1&-1-2i,\ -1+6i\\
&2&-5+2i\\
&3&-3,\ 3-2i,\ 1+4i,\ 1-4i,\ -5-2i\\
+7+10i&0&1+4i,\ 1-4i,\ -1+6i,\ -1-6i\\
&1&-1+2i,\ 3+2i,\ -5+2i\\
&2&1+i,\ 3-2i\\
&3&-1-2i,\ -3,\ -5-2i\\
+11+6i&0&1+i,\ -1+2i,\ -3,\ 1+4i,\ 1-4i,\ -7\\
&1&-1-2i,\ 3+2i,\ 3-2i\\
&2&-5-2i,\ -1+6i,\ 5-4i\\
&3&-5+2i,\ 5+4i,\ 7-2i.
\end{array}
\]

\Article{63}
Wir werden nun versuchen, durch die Untersuchung von Zahlenbeispielen gemeinsame Kriterien für die Moduln zu finden, für die eine gegebene Primzahl denselben Charakter hat. Die Moduln setzen wir stets als primär unter ihren Assoziierten voraus, also als solche Zahlen $a+bi$, für die entweder $a\equiv1,b\equiv0$ oder $a\equiv3,b\equiv2\pmod4$ gilt.

\Needspace{15\baselineskip}
Für die Zahl $1+i$, mit der wir beginnen, lässt sich die in den Beispielen erkennbare Regel leichter erfassen, wenn wir die Moduln der ersten Art, für die $a\equiv1,b\equiv0$ gilt, von den Moduln der zweiten Art, für die $a\equiv3,b\equiv2$ gilt, trennen. Mithilfe der Tabelle des vorigen Artikels finden wir die folgenden Zuordnungen:
\SourceLeaf{145}%
\[
\begin{array}{c|l}
\text{Charakter}&\text{Moduln der ersten Art.}\\\hline
0&5+4i,\ -7+8i,\ -7-8i,\ -11+4i\\
1&1-4i,\ -3+8i,\ -3-8i,\ 9+4i,\ -11\\
2&5-4i,\ -7,\ -11-4i\\
3&-3,\ 1+4i,\ 5+8i,\ 5-8i,\ 9-4i
\end{array}
\]
Betrachten wir diese siebzehn Beispiele aufmerksam, so finden wir in allen den Charakter $\equiv\frac14(a-b-1)\pmod4$.

Ebenso gelten die Zuordnungen
\[
\begin{array}{c|l}
\text{Charakter}&\text{Moduln der zweiten Art.}\\\hline
0&3-2i,\ -1-6i,\ 7+2i,\ -5+6i,\ -1+10i,\ 11+6i\\
1&-5+2i,\ -1+6i,\ 7-2i,\ -1-10i,\ 3+10i\\
2&-1+2i,\ -5-2i,\ 3-10i,\ 7+10i\\
3&-1-2i,\ 3+2i,\ -5-6i,\ 7-10i,\ 11-6i
\end{array}
\]
In allen diesen zwanzig Beispielen findet man bei einiger Aufmerksamkeit den Charakter $\equiv\frac14(a-b-5)\pmod4$.

Diese beiden Regeln lassen sich leicht zu einer für beide Arten von Moduln gültigen Regel zusammenfassen, wenn wir bedenken, dass $\frac14bb$ für Moduln der ersten Art $\equiv0$, für Moduln der zweiten Art dagegen $\equiv1\pmod4$ ist. Somit ist der Charakter der Zahl $1+i$ bezüglich jedes primen, unter seinen Assoziierten primären Moduls $\equiv\frac14(a-b-1-bb)\pmod4$.

Nebenbei ist zu bemerken, dass, da $(b+1)^2$ stets die Form $8n+1$ hat oder gleichbedeutend $\frac14(2b+bb)$ gerade ist, jener Charakter stets gerade oder ungerade ist, je nachdem, ob $\frac14(a+b-1)$ gerade oder ungerade ist. Das stimmt mit der in Artikel~58 angegebenen Regel für den quadratischen Charakter überein.

Da $\frac14(a-b-1),\frac14(a-b+3)$ ganze Zahlen sind, von denen eine gerade, die andere ungerade ist, ist ihr Produkt gerade, oder gleichbedeutend $\frac18(a-b-1)(a-b+3)\equiv0\pmod4$. Daher kann anstelle des angegebenen Ausdrucks für den biquadratischen Charakter auch der folgende genommen werden:
\[
\frac14(a-b-1-bb)-\frac18(a-b-1)(a-b+3)=\frac18(-aa+2ab-3bb+1)
\]
\Needspace{7\baselineskip}
Diese Form empfiehlt sich auch dadurch, dass sie nicht auf primäre Moduln beschränkt ist, sondern nur voraussetzt, dass $a$ ungerade und $b$ gerade ist. Denn unter dieser Voraussetzung ist offenbar entweder $a+bi$ oder $-a-bi$ unter seinen Assoziierten primär, und der Wert jener Formel ist für beide Moduln derselbe.
